\documentclass[11pt]{article}
\usepackage[margin=25mm]{geometry}
\usepackage{fontspec}
\setmainfont{Latin Modern Roman}
\usepackage{amsmath,amssymb,mathtools}
\usepackage{unicode-math}
\setmathfont{Latin Modern Math}
\usepackage{xurl,hyperref}
\hypersetup{colorlinks=true,urlcolor=blue}
\setcounter{secnumdepth}{0}
\setlength{\parindent}{0pt}
\setlength{\parskip}{5pt}
\setlength{\emergencystretch}{3em}
\newcommand{\sref}[2]{\hyperlink{source:#1:#2}{[S#2]}}
\newcommand{\cataloguescope}[1]{\paragraph{Recorded target.}#1}
\newcommand{\R}{\mathbb R}
\newcommand{\C}{\mathbb C}
\newcommand{\N}{\mathbb N}
\newcommand{\Z}{\mathbb Z}
\newcommand{\Q}{\mathbb Q}
\newcommand{\F}{\mathbb F}
\newcommand{\eps}{\varepsilon}
\begin{document}
\section{Hoop Conjecture for Black-Hole Formation}\label{statement}

\subsection{Definitions and mathematical statement}

\paragraph{Shared notation.}$\mathbb N=\{1,2,\ldots\}$; the standard number fields have their usual meanings. All additional conventions are stated locally.


\paragraph{Definitions and conventions.} Use geometrized units $G=c=1$, so a mass $M$ has units of length. An event horizon is the boundary separating points that can send causal signals to future null infinity from the black-hole region. The hoop conjecture proposes a relation between compactness, a notion of mass, and circumferences in all directions. In time-symmetric initial data an apparent horizon corresponds to an appropriate closed minimal surface, which is a more specific geometric setting discussed in the report.

\paragraph{Source scope and qualifications.} The primary report explicitly says that mass and circumference are not clearly defined in this formulation. Event and apparent horizons are different notions. No arbitrary definition of a hoop, quasi-local mass, or horizon is imposed here, and the time-symmetric results do not resolve every interpretation.

\paragraph{Catalogue mathematical target.}
\paragraph{Catalogue wording retained for traceability.}
Do black holes with horizons form if and only if a mass $M$ is compacted into a region whose circumference in every direction is at most $4\pi M$?

\paragraph{Saved literature context.} The supplied catalogue's assessment dated 2026-08-21 records: The extracted item is explicitly labelled a loosely formulated conjecture. Without the omitted mathematical formulation there is no verified theorem or counterexample to classify it as solved or disproved. The assessment is retained as dated provenance; it is not a new comprehensive literature or open-status certification.
\subsection{Short English statement}
Make precise and test the claim that horizons form exactly when mass is compressed inside sufficiently short hoops in every direction.
\subsection{Research attempt}
\paragraph{Model, date and investigated scope.}
OpenAI GPT-6 Astra (configured model \texttt{gpt-6-astra}), ultra reasoning; generated 11 October 2026. Reconstructs the constant 4πM at the Schwarzschild throat and exposes definitional obstacles beyond spherical symmetry.

\paragraph{Sources and literature check.}
The complete supplied definition, including source qualifications and the dated literature assessment, was reviewed. The calculation below is an elementary restricted attempt motivated by [S2]; it does not claim novelty or a fresh literature-status determination.

\paragraph{Context and investigated scope.} The complete supplied definition, including source qualifications and the dated literature assessment, was reviewed. The calculation below is an elementary restricted attempt motivated by [S2]; it does not claim novelty or a fresh literature-status determination.

\paragraph{The spherical equality check.} On the time-symmetric Schwarzschild exterior of mass $M>0$, use areal radius $r>2M$ and metric
\[
 g=(1-2M/r)^{-1}dr^2+r^2d\Omega^2.
\]
The outward unit normal to a coordinate sphere is $\nu=\sqrt{1-2M/r}\,\partial_r$. Its area is $4\pi r^2$, so its mean curvature, the normal logarithmic area derivative, is
\[
 H=\nu\bigl(\log(4\pi r^2)\bigr)=\frac2r\sqrt{1-2M/r}.
\]
It tends to zero at $r=2M$. In the smooth doubled spatial geometry this limiting sphere is the throat minimal surface. Every great circle there has length $2\pi r=4\pi M$, reproducing the proposed constant in this symmetric apparent-horizon example.

\paragraph{Why circumference alone does not generalize.} The calculation uses the special facts that all great-circle circumferences agree and that $r$ is geometrically fixed by area. Away from symmetry one cannot replace this by an arbitrary coordinate circumference. Even a fixed round sphere has closed curves of arbitrarily large length by repeated winding, so ``every direction'' cannot mean ``every closed curve''. Nor does the local equation $H=0$ identify an event horizon, which uses the entire spacetime future.

\paragraph{Remaining gap.} This is a consistency test of the normalization, not an if-and-only-if formation theorem. The first unresolved issue is the source's own missing definitions: which quasi-local mass, which family of hoops, and which horizon notion are intended. After fixing them, both necessity and sufficiency for nonspherical dynamical data still require proof. At $M=0$ the displayed exterior becomes Euclidean space and has no positive-radius coordinate minimal sphere, consistent with the degenerate limit.

\paragraph{Verification and adversarial audit.}
Checked unit normal and mean curvature, the doubled-throat interpretation, zero mass, and event/apparent-horizon distinction. Model self-audit only; no independent human proof checking. Prompts A and B were applied to this canonical record and its qualified source scope. Checked unit normal and mean curvature, the doubled-throat interpretation, zero mass, and event/apparent-horizon distinction. Retained unresolved status and zero original-target progress; A precise mass/hoop/horizon formulation is needed before either general implication can be proved.

Coordinator checked the Schwarzschild mean curvature and round-horizon circumference, including the failure of unrestricted curve lengths to define a hoop criterion. This additional check was model-assisted, not independent human verification.

\paragraph{Tools, process and contributions.}
Python 3 was used for file preparation and structural checks; no numerical experiment or theorem prover was used. Contribution prompt A was followed by prompt B. GPT-6 Astra with ultra reasoning drafted the research and performed its stated self-audit; any additional coordinator contribution is identified in the verification and audit. A GPT-6-based Codex coordinator assembled the manuscript and checked repository integration; its exact variant and reasoning configuration were not exposed. No human mathematical contribution, independent human proof check or formal proof certification is supplied. Complete derivations and the consulted source references are included here; private reasoning transcripts are not published.

\paragraph{Final status and first remaining gap.}
\textbf{UNRESOLVED}. A precise mass/hoop/horizon formulation is needed before either general implication can be proved.

\section{Completion Estimate}
\noindent\textbf{COMPLETION: 0\%}
The calculation is an auxiliary or restricted attempt and establishes no newly verified portion of the full source target. A precise mass/hoop/horizon formulation is needed before either general implication can be proved.

\subsection{Sources}
\begin{itemize}
\item[\hypertarget{source:statement:1}{S1}] ULAM.AI and the UnsolvedMath Contributors, UnsolvedMath, catalogue entry. Catalogue curation under CC BY 4.0; definitions and English exposition added locally.
\url{https://huggingface.co/datasets/ulamai/UnsolvedMath}.
\item[\hypertarget{source:statement:2}{S2}] Mathematical Aspects of General Relativity (2006), Oberwolfach Reports, DOI 10.4171/owr/2006/02, pp. 73--118; official report source and mathematical context.
\url{https://doi.org/10.4171/owr/2006/02}.
\item[\hypertarget{source:statement:3}{S3}] Official EMS Press PDF of the same Oberwolfach report; PDF pages 16--17, printed pages 88--89.
\url{https://ems.press/content/serial-article-files/46032}.
\end{itemize}
\label{end:statement}
\end{document}
